%%%%%%%%%% %%%%%%%%%% %%%%%%%%%% %%%%%%%%%% 
\section{Demystifying RAG Computation}
\label{sec:02BGnMotivation}
% Text matter begin
%%%%%%%%%% %%%%%%%%%% %%%%%%%%%% %%%%%%%%%% 
%%%%%%%%%% 
%% \documentclass{article}
% \usepackage[table]{xcolor} % For cell coloring
% \usepackage{makecell} % For multi-line cells
% \usepackage{pifont} % For tick and cross marks
% \usepackage{graphicx} % For rotating text
% \usepackage{booktabs} % For better table styling

% \newcommand{\cmark}{\ding{51}} % Tick mark (✓)
% \newcommand{\xmark}{\ding{55}} % Cross mark (✗)



\begin{table}[]
\centering
\small
\renewcommand{\arraystretch}{1.2} % Adjust row height
\setlength{\tabcolsep}{5pt} % Adjust column spacing

\begin{tabular}{|c|c|c|c|c|}
\hline
\textbf{Features} & 
% \rotatebox{90}{\textbf{FlashRAG~\cite{jin2024flashrag}}} & 
% \rotatebox{90}{\textbf{EdgeRAG~\cite{EdgeRAG}}} & 
% \rotatebox{90}{\textbf{PipeRAG~\cite{jiang2024piperag}}} & 
% \rotatebox{90}{\textbf{MaestroRAG}} \\\hline
\rotatebox{0}{\textbf{FR}} & 
\rotatebox{0}{\textbf{ER}} & 
\rotatebox{0}{\textbf{PR}} & 
\rotatebox{0}{\textbf{MR}} \\\hline

\textbf{Heterogeneous platform usage} & \cellcolor{green!20} \cmark & \cellcolor{green!20} \cmark & \cellcolor{green!20} \cmark & \cellcolor{green!20} \cmark \\\hline
\textbf{Caching in retrieval} & \cellcolor{red!20} \xmark & \cellcolor{green!20} \cmark & \cellcolor{red!20} \xmark & \cellcolor{green!20} \cmark \\\hline
\textbf{Pipelining CPU-GPU} & \cellcolor{red!20} \xmark & \cellcolor{red!20} \xmark & \cellcolor{green!20} \cmark & \cellcolor{green!20} \cmark \\\hline
\textbf{Concurrent encoding/retrieval} & \cellcolor{red!20} \xmark & \cellcolor{red!20} \xmark & \cellcolor{red!20} \xmark & \cellcolor{green!20} \cmark \\\hline
\textbf{Multiple CPU workers} & \cellcolor{red!20} \xmark & \cellcolor{red!20} \xmark & \cellcolor{red!20} \xmark & \cellcolor{green!20} \cmark \\\hline
\textbf{Async worker execution} & \cellcolor{red!20} \xmark & \cellcolor{red!20} \xmark & \cellcolor{red!20} \xmark & \cellcolor{green!20} \cmark \\\hline
\textbf{Input-sensitive caching} & \cellcolor{red!20} \xmark & \cellcolor{red!20} \xmark & \cellcolor{red!20} \xmark & \cellcolor{green!20} \cmark \\\hline
\end{tabular}
\caption{Comparison of \design{} with state-of-the-art resource management techniques. FR=\flashRAG{}, ER=\edgeRAG{}, PR=\pipeRAG, MR=\design{}}
\label{tab:rag_comparison}
%\vspace*{-0.3in}
\end{table}


% \usepackage[table]{xcolor} % For cell coloring
% \usepackage{makecell} % For multi-line cells
% \usepackage{pifont} % For tick and cross marks

% \newcommand{\cmark}{\ding{51}} % Tick mark
% \newcommand{\xmark}{\ding{55}} % Cross mark

%

% \begin{table}[!ht]
% \caption{Comparison of our work with state-of-the-art resource management approaches.}
% \centering
% \scriptsize
% \renewcommand{\arraystretch}{1.2} % Adjust row height
% \setlength{\tabcolsep}{5pt} % Adjust column spacing

% \begin{tabular}{|p{4.2cm}|c|c|c|c|}
% \hline
% \textbf{List of Optimizations} & \textbf{FlashRAG} & \textbf{EdgeRAG} & \textbf{PipeRAG} & \textbf{Ours} \\\hline
% \textbf{Heterogeneous platform usage} & \cellcolor{green!20} \cmark & \cellcolor{green!20} \cmark & \cellcolor{green!20} \cmark & \cellcolor{green!20} \cmark \\\hline
% \textbf{Caching in retrieval} & \cellcolor{red!20} \xmark & \cellcolor{green!20} \cmark & \cellcolor{red!20} \xmark & \cellcolor{green!20} \cmark \\\hline
% \textbf{Pipelining CPU and GPU} & \cellcolor{red!20} \xmark & \cellcolor{red!20} \xmark & \cellcolor{green!20} \cmark & \cellcolor{green!20} \cmark \\\hline
% \makecell{\textbf{Concurrent execution of} \\ \textbf{encoding and retrieval}} & \cellcolor{red!20} \xmark & \cellcolor{red!20} \xmark & \cellcolor{red!20} \xmark & \cellcolor{green!20} \cmark \\\hline
% \textbf{Multiple workers on CPU} & \cellcolor{red!20} \xmark & \cellcolor{red!20} \xmark & \cellcolor{red!20} \xmark & \cellcolor{green!20} \cmark \\\hline
% \textbf{Async worker execution} & \cellcolor{red!20} \xmark & \cellcolor{red!20} \xmark & \cellcolor{red!20} \xmark & \cellcolor{green!20} \cmark \\\hline
% \textbf{Input-sensitive caching} & \cellcolor{red!20} \xmark & \cellcolor{red!20} \xmark & \cellcolor{red!20} \xmark & \cellcolor{green!20} \cmark \\\hline
% \end{tabular}
% \label{tab:rag_comparison}
% \end{table}

% \usepackage[table]{xcolor} % For cell coloring
% \usepackage{makecell} % For multi-line cells
% \usepackage{pifont} % For tick and cross marks
% \usepackage{graphicx} % For rotating text

% \newcommand{\cmark}{\ding{51}} % Tick mark
% \newcommand{\xmark}{\ding{55}} % Cross mark

% \begin{document}

% \begin{table}[!ht]
% \caption{Comparison of our approach with state-of-the-art resource management techniques.}
% \centering
% \small
% \renewcommand{\arraystretch}{1.2} % Adjust row height
% \setlength{\tabcolsep}{3pt} % Reduce column spacing

% \begin{tabular}{|p{6cm}|c|c|c|c|}
% \hline
% \textbf{Optimization} & 
% \rotatebox{90}{\textbf{FlashRAG~\cite{jin2024flashrag}}} & 
% \rotatebox{90}{\textbf{EdgeRAG~\cite{EdgeRAG}}} & 
% \rotatebox{90}{\textbf{PipeRAG\cite{jiang2024piperag}}} & 
% \rotatebox{90}{\textbf{Ours}} \\\hline

% Heterogeneous platform usage & \cellcolor{green!20} \cmark & \cellcolor{green!20} \cmark & \cellcolor{green!20} \cmark & \cellcolor{green!20} \cmark \\\hline
% Caching in retrieval & \cellcolor{red!20} \xmark & \cellcolor{green!20} \cmark & \cellcolor{red!20} \xmark & \cellcolor{green!20} \cmark \\\hline
% Pipelining CPU-GPU & \cellcolor{red!20} \xmark & \cellcolor{red!20} \xmark & \cellcolor{green!20} \cmark & \cellcolor{green!20} \cmark \\\hline
% Concurrent execution (encoding \& retrieval) & \cellcolor{red!20} \xmark & \cellcolor{red!20} \xmark & \cellcolor{red!20} \xmark & \cellcolor{green!20} \cmark \\\hline
% Multiple CPU workers & \cellcolor{red!20} \xmark & \cellcolor{red!20} \xmark & \cellcolor{red!20} \xmark & \cellcolor{green!20} \cmark \\\hline
% Async worker execution & \cellcolor{red!20} \xmark & \cellcolor{red!20} \xmark & \cellcolor{red!20} \xmark & \cellcolor{green!20} \cmark \\\hline
% Input-sensitive caching & \cellcolor{red!20} \xmark & \cellcolor{red!20} \xmark & \cellcolor{red!20} \xmark & \cellcolor{green!20} \cmark \\\hline
% \end{tabular}
% \label{tab:rag_comparison}
% \end{table}




%%%%%%%%%% 

The RAG pipeline enables LLMs to deliver enhanced personalization, use-case specificity, and customization, while preserving user data security. At its core, a compact generative model produces responses augmented with a set of \topk{} relevant knowledge items. As shown in \Cref{fig:RAG_background}, the pipeline begins by encoding the user query into high-dimensional embeddings. A retriever then performs a similarity search over the user’s local database to extract the \topk{} pertinent entries, which are combined with the original query and passed to the LLM to generate the final response. We next discuss the retrieval, augmentation, and generation stages in detail.

\subsection{Different Stages of a RAG System}
\label{sec:BG:RAGstage}
%\subsection{Encoding}
\noindent\paragraph{Encoding}
\label{sec:BG:Encoding}
A RAG pipeline begins with \emph{encoding}, where the natural-language query is tokenized into subword units and passed through a neural encoder (e.g., \texttt{BERT}~\cite{devlin2019bert}, \texttt{RoBERTa}~\cite{liu2019roberta}, \texttt{T5}~\cite{raffel2020exploring}, or compact variants like \texttt{DistilBERT}~\cite{sanh2019distilbert}, \texttt{e5-base-v2}~\cite{wang2022text}), to generate high-dimensional embeddings. These models, despite varying in size and efficiency, map text into a semantic vector space. The forward pass, dominated by multi-head attention and fully connected layers, constitutes most of the encoder’s runtime. The resulting embeddings are then used in the retrieval stage to match against candidate knowledge elements.

%\subsection{Retrieval}
\noindent\paragraph{Retrieval}
\label{sec:BG:Retrieval}
Following encoding, the system searches a \emph{vector database} (constructed by encoding each knowledge item with the same model)to find relevant entries. This shared embedding space enables direct similarity comparisons, while database construction, being computationally intensive, is usually performed offline or during idle periods.

For efficient lookup, the system builds either hierarchical (tree-based) indices with sublinear search but higher memory cost or flat indices that perform exhaustive searches with higher time complexity. Retrieval computes a similarity metric (e.g., cosine or L2 distance~\cite{l2}) between the query embedding and database entries to return \topk{} matches. Depending on resources, this computation can run on CPUs or accelerators, though CPUs often excel for large, I/O-intensive indices. Finally, the top entries are passed to the \emph{augmentation} stage for prompt composition and generation.

%\subsection{Augmentation}
\noindent\paragraph{Augmentation}
\label{sec:BG:Augmentation}
After retrieval, the augmentation stage constructs a contextualized prompt for the generation model by mapping each of the \topk{} retrieved embeddings to their original text and merging these excerpts with the user’s query. Optional style or verbosity preferences (e.g., concise, detailed, formal) may also be appended. This augmented prompt enables the generation phase to produce accurate and personalized outputs using the retrieved knowledge and user instructions.

%\subsection{Generation}
\noindent\paragraph{Generation}
\label{sec:BG:Generation}
The final stage of the RAG pipeline generates natural-language outputs using off-the-shelf LLMs (e.g., Llama 8B~\cite{llama_3_1_8b_instruct}). Because this stage involves parallel operations such as GEMM and self-attention, GPUs are typically used for efficient execution~\cite{attention_gpu}. However, in edge devices, GPUs have limited VRAM and share resources with other tasks, so generative models are stored on device storage and loaded into GPU memory only when needed, preventing VRAM contention and improving resource utilization. 

\subsection{System Considerations}
\label{sec:classical_pipeline}
The fundamental RAG stages—encoding, retrieval, augmentation, and generation—are consistent across datacenter servers, desktops, and edge devices, but their execution varies widely with hardware constraints.

A typical pipeline loads the encoder model from storage into CPU memory, a significant overhead for large models under tight memory budgets. Servers usually keep both encoder and generative models resident in memory to support large batches~\cite{jiang2024piperag, jin2024flashrag, jiang2024megascale}, whereas consumer grade devices may rely on swapping,
%desktops may swap model segments, and edge devices rely on 
smaller models or aggressive offloading~\cite{EdgeRAG,NanoLLM}. Once loaded, encoding often uses GPUs for parallelism, though multicore CPUs can suffice for smaller models or workloads. CPU–GPU data transfers should be carefully managed to balance throughput against transfer costs. Retrieval, which loads an index in CPU memory to locate the \topk{} matches, is typically I/O- and memory-bound, with costs rising as the database grows or updates frequently. Augmentation, in contrast, is lightweight and merges retrieved items with the user’s query and style preferences before sending the prompt to the LLM (usually on a GPU). Edge platforms with unified memory~\cite{ying2022exploiting} can simplify pointer sharing but risk contention when processing large inputs~\cite{kim2020batch, dashti2017analyzing, cavicchioli2017memory, jeong2012qos}.

Batching is used to amortize overhead and improve throughput, but large batches can inflate memory use and cause thermal throttling~\cite{na2021scalable, stojkovic2025tapas}. Concurrent execution across stages can reduce latency but risks pipeline back-pressure or idle cycles if stages are imbalanced. Overall, efficient RAG execution requires careful coordination of CPU- and GPU-bound stages, I/O- and memory-heavy operations, and hardware power limits. \Cref{sec:03Characterization} explores these trade-offs in greater detail.

% \subsection{State-of-the-art RAG Systems}  
% \label{sec:rag-sota}
% \red{maybe merge with intro SOTA. I tried in intro. Let us remove this section after editing the intro and related work sections.  }
% %Modern RAG pipelines typically perform encoding and retrieval on either CPUs or GPUs, followed by CPU-based augmentation and GPU-based generation~\cite{jin2024flashrag, jiang2024piperag, EdgeRAG}. 
% FlashRAG~\cite{jin2024flashrag} emphasizes GPU-based encoding and retrieval, which works well in large-scale deployments but is less effective on edge devices with lower request rates. 
% \pipeRAG{} runs CPU-based encoding and retrieval, then batches results for one-shot GPU generation, improving throughput but leaving pipeline bubbles when CPU stages do not fully overlap with generation; frequent model and index reloads also incur I/O overhead. 
% \edgeRAG{} reduces memory use by caching coarse indices and generating fine-grained embeddings on demand, but this increases compute and power costs, which are critical for edge devices, and sacrifices accuracy due to its reliance on tree-based IVF indices (\ref{sec:BG:Retrieval}) over flat indices. 
%%%%%%%%%% %%%%%%%%%% %%%%%%%%%% %%%%%%%%%% 


%%%%%%%%%% %%%%%%%%%% %%%%%%%%%% %%%%%%%%%% 
\section{Detailed Characterization}
\label{sec:03Characterization}

%%%%%%%%%% %%%%%%%%%% %%%%%%%%%% %%%%%%%%%% 
This section presents a comprehensive performance characterization of a RAG pipeline on a desktop-grade machine, focusing on how each stage performs under different resource allocations, database sizes, and batching strategies. Our goal is to gain insight into the computational and I/O demands of the various RAG stages, thus motivating the pipeline design described below. All characterization experiments are performed on a machine consisting of an Intel i9-14900K~\cite{intel_cpu} processor (24 physical cores), 128\,GB system memory and RTX 4090 GPU~\cite{nvidia_rtx_gpu}. \emph{We limit generation lengths to 120 tokens}. To better explain the end-to-end pipeline, we decompose the RAG workflow into four stages: \emph{encode}, \emph{retrieve}, \emph{augment}, and \emph{generate} due to their divergent computational and memory footprints (discussed in more detail in \Cref{sec:04Design}). The first three stages are run on the CPU, while the last stage is performed on the GPU. \Cref{Fig:stacked_graph} shows the execution time breakdown of a RAG request, where the encoding of prompts, the retrieval of knowledge elements, and the augmentation stages are performed on a single CPU core (to get the compute ratio of each stage). This experiment is done for a batch size of 8 queries retrieving knowledge elements from a 4M size of database. Here, and henceforth in this paper, by ``4M size of database'' we mean that there are 4M knowledge elements, where each knowledge element is a group of $\approx100$ words.

\subsection{Encoding Stage}
\label{sec:char:Encoding}
As discussed in \ref{sec:BG:Encoding}, the encoding stage tokenizes the input and performs a projection, converting them into a higher dimensional embeddings.
\Cref{Fig:forwardpass_8cores} shows the overall latency for a forward pass with increased batch size (\texttt{BS}) for the encode operation over a fixed number of CPU cores (8 performance cores). 
We find that the overall latency of this stage is a function of the computational resources (\#cores) assigned to it and the amount of work given to it (\texttt{BS}). The latency reduction by assigning a higher number of cores quickly decreases and saturates as we increase the core count from 8 to 16, as shown in \Cref{Fig:forwardpass_bs16_range_discrete}. Thus, provisioning more cores for encoding stage yields diminishing returns unless amortized by a commensurate number of parallel requests (e.g., \texttt{BS=32}). 

\subsection{Retrieval stage}
\label{sec:char:Retrieval}
As discussed in \ref{sec:BG:Retrieval}, the retrieval stage identifies the \topk most similar entries from a stored vector database through embedding similarity search. This stage is highly I/O and memory bound, requiring intricate characterization to understand bottlenecks to make informed design decisions. 
The retrieval stage execution is comprised of two parts: i) index fetch and ii) similarity search. Fetching of indices occurs once for a batch and is independent of the batch size, while the similarity search happens for each sample in a batch. \Cref{Fig:batchsize_stacked_linear_labeled} plots the contribution of index fetch and similarity search in retrieval stage execution. As expected, index fetch incurs a fixed latency for a fixed database size, while the similarity search latency increases with batch size. We also note that index fetch dominates the execution time at low batch sizes, while similarity search dominates at higher batch sizes. When the batch size is 32, the total execution time increases considerably due to the saturation of the DRAM bandwidth.

\Cref{Fig:latency_cores} highlights the multicore scaling in the retrieval stage, where we observe that for a fixed number of requests in a batch, with increasing CPU cores, the latency gradually decreases. The scaling is sub-linear in nature, indicating diminishing returns on increase in computational investment. However, the trend of latency with respect to database size is more linear as shown in \Cref{Fig:latency_dbsize}. These results suggest that latency is more strongly correlated with IO processes than with computational jobs in the retrieval phase. 

\noindent\textbf{Execution Breakdown Across Stages: }
\Cref{Fig:latency_cores} highlights how the overall CPU latency drops as we increase the number of cores for retrieval. Although doubling the cores from 8 to 16 continues to lower latency, the gains taper off due to memory and I/O constraints. 
When examining different database sizes in \Cref{Fig:latency_dbsize}, we observe near-linear increases in retrieval time as the database grows from 1\,M to 8\,M. Large indices incur more frequent cache misses and potential page faults, increasing latencies. 
Meanwhile, \Cref{Fig:batchsize_stacked_linear_labeled} shows a deeper breakdown of the retrieval time when varying batch sizes: as the batch size grows, the time attributed to the I/O operations (labeled ``Index Fetch'') becomes a significant portion of total retrieval, while the similarity search component may remain modest unless the batch is large enough to saturate the CPU compute.

\noindent\textbf{Impact of Batch Size and Multicore Scaling:} 
~\Cref{Fig:forwardpass_8cores} shows how encoding time scales with batch size on 8 CPU cores. Although small batch sizes observe lower latency compared to larger batch sizes, they do not fully utilize CPU resources. However, larger batches risk higher memory footprints and potential thrashing. In retrieval, these tradeoffs manifest in \Cref{Fig:batchsize_stacked_linear_labeled}, where going beyond a batch size of 16 significantly inflates total execution time, partly because the system expends more time on either index fetching or computing similarities for many queries in one go. 

\noindent\textbf{Effects of Database Size:}
\Cref{Fig:latency_dbsize} demonstrates how the retrieval latency scales with the size of the DB on 16 CPU cores. Increasing from 1\,M to 2\,M entries roughly doubles the time, and at 8\,M entries, we see further linear growth. While more CPU cores can mitigate retrieval cost to a degree, the cost of scanning a large index remains partly I/O-bound. As a result, simply adding cores does not proportionally reduce the penalty incurred by a growing database.

\subsection{Key Observations and Motivation for Design}
We observe that encoding scales well with moderate core counts and remains mostly compute-bound, while retrieval is especially sensitive to memory footprint and I/O activity, causing diminishing returns as the batch size grows. Augmentation is lightweight and best merged with retrieval to avoid extra data transfers. Although GPU-based generation can be the longest stage, its runtime often remains stable unless the model is repeatedly reloaded or the prompts are excessively large. Consequently, it becomes crucial to allocate CPU cores judiciously for encoding and retrieval and to prevent ballooning memory usage at high batch sizes or large databases. Furthermore, offloading the encoder to the GPU can be inefficient when transfers dominate, emphasizing the importance of selectively assigning each stage to CPU or GPU. These insights form the basis for our multi-stage pipeline and scheduling in \Cref{sec:04Design}.
%%%%%%%%%% %%%%%%%%%% %%%%%%%%%% %%%%%%%%%%
% Tex matter end